%%%%%%%%%%%%%%%%%%%%%%%%%%%%%%%%%%%%%%%%%
% Specjalna strona pracy ze streszczeniem i abstractem w j. angielskim
% Szablon pracy dyplomowej
% Wydział Informatyki 
% Zachodniopomorski Uniwersytet Technologiczny w Szczecinie
% autor Joanna Kołodziejczyk (jkolodziejczyk@zut.edu.pl)
% Bardzo wczesnym pierwowzorem szablonu był
% The Legrand Orange Book
% Version 5.0 (29/05/2025)
%
% Modifications to LOB assigned by %JK
%%%%%%%%%%%%%%%%%%%%%%%%%%%%%%%%%%%%%%%%%


\begin{center}
\noindent {{\color{blueZUT}\Large\sffamily  {Streszczenie}}}\\[1cm] 
\end{center}

Celem pracy jest opracowanie i implementacja aplikacji umożliwiającej bezdotykowa sterowanie kursorem myszy za pomocą sygnałów SSVEP (Steady-State Visual Evoked Potentials). Aplikacja pozwoli na kontrolę ruchu kursora w czterech kierunkach (góra, dół, lewo, prawo) oraz symulowanie kliknięcia lewym przyciskiem myszy. Użytkownik będzie sterował kursorem poprzez skupienie wzroku na wyświetlanych na ekranie polach migających z ustalonymi częstotliwościami. Aplikacja będzie składała się z trzech głównych modułów, obejmujących: moduł do odczytu i przetworzenia sygnału EEG; stymulator ekranowy; oraz moduł sterujący pracą myszki. Praca zostanie wykonana w środowisku Python z wykorzystaniem specjalizowanej biblioteki do odczytu przetwarzania sygnału EEG (mne).

\vspace{10pt}
\noindent{\bf słowa kluczowe:} SSVEP, EEG, interfejs mózg-komputer (BCI), Python

\vfill

\begin{center}
\noindent {{\color{blueZUT}\Large\sffamily {Abstract}}}\\[1cm] 
\end{center}
The aim of this work is to develop and implement an application enabling touchless mouse cursor control using SSVEP (Steady-State Visual Evoked Potentials) signals. The application will allow for cursor movement control in four directions (up, down, left, right) and simulate left-clicks. The user will control the cursor by focusing their gaze on on-screen fields that flash at predetermined frequencies. The application will consist of three main modules: an EEG signal acquisition and processing module; an on-screen stimulator; and a mouse control module. The work will be implemented in Python using a specialized EEG signal acquisition and processing library (mne).


\vspace{10pt}
\noindent{\bf keywords:} SSVEP, EEG, mouse cursor control, brain-computer interface (BCI), Python